\chapter*{Abstract}
\addcontentsline{toc}{chapter}{Abstract}
\begin{center}\large\bfseries\thesistitle\end{center}

Surface finishing often requires local rework around a selected region rather than uniform treatment of an entire component. A local polishing operation must coordinate tool movement, contact force, and processing time within that region. In the setting considered here, a person specifies the region to be treated, avoiding dependence on autonomous recognition of fine scratches with the available equipment. This spatial instruction identifies where processing should occur but leaves the motion--force sequence to be constructed. Rules and demonstration replay provide legitimate solutions; the research question is when observation-conditioned skill generation provides additional execution value.

This thesis develops a human-guided local polishing framework that connects region specification with temporally paired motion--force demonstrations. A tracked teaching system records tool pose, three-axis interaction force, and RGB observations. Task-centered position learning is paired with an implemented direction-alignment adapter that maps observations into the teaching frame and restores generated XY paths. A conditional Flow Matching policy generates sequences of position, orientation, and target force from images, robot state, and measured-force history. A shared execution layer restores task coordinates, connects plans, and conditions commands before an existing lower-level force controller regulates contact. Skill generation, contact regulation, and procedural completion evaluation have distinct responsibilities.

Two robot evaluations examine this division, with direction alignment disabled. E1 compares rule-based planning, force-inclusive demonstration replay, and the force-conditioned policy for a fixed nominal task. All five executions per method complete the required contact--loaded-pass--exit procedure. Replay has a shorter mean completion time and smaller time dispersion than the learned policy, whose excess-force exposure is greater. E2 compares motion-only learning and joint motion--force learning with and without measured-force observation. The force-observing policy has a \BCPercent\% shorter mean operator-based completion interval than its force-masked counterpart in the recorded cohort. One motion-only attempt remains in contact and is interrupted. These results establish an executable, human-guided motion--force skill and reveal differences in timing, phase progression, and force use across configurations. Rules and replay remain effective alternatives for the fixed task. The evaluation concerns procedural execution rather than measured surface-quality improvement or robustness across changed task conditions.

\vspace{0.4cm}
\noindent\textbf{Keywords:} Local Robotic Polishing, Human Task Specification, Motion--Force Skill, Imitation Learning, Flow Matching, Contact-Task Evaluation
